% problem-000034 4 高压钠氯及钠氦晶体
\section*{4. 高压钠氯及钠氦晶体}
\textit{下列为分问。}

\subsection*{4-1}
NaCl晶体在50\~300GPa的⾼压下和Na或Cl 反应，可以形成不同组成、不同结构的晶体。左图给出其中三种晶体的晶胞（⼤球为氯原⼦，⼩球为钠原⼦）。写出A、B、C的化学式。

（图：）
A

（图：）
B

（图：）
C

（图：）

\subsection*{4-2}
在超⾼压(300GPa)下，⾦属钠和氦可形成化合物。结构中，钠离⼦按简单⽴⽅排布，形成 $\mathrm { N a } _ { 8 }$ ⽴⽅体空隙（如右图所⽰），电⼦对(2e−)和氦原⼦交替分布填充在⽴⽅体的中⼼。

\subsection*{4-2}
-1 写出晶胞中的钠离⼦数。

\subsection*{4-2}
-2 写出体现该化合物结构特点的化学式。

\subsection*{4-2}
-3 若将氦原⼦放在晶胞顶点，写出所有电⼦对 $( 2 \mathrm { e ^ { - } } )$ 在晶胞中的位置。

\subsection*{4-2}
-4 晶胞边⻓ $a = 3 9 5 \mathrm { p m }$ ，计算此结构中 $_ \mathrm { N a - H e }$ 的间距�和晶体的密度�（单位： $\mathrm { g c m ^ { - 3 } } )$
